\documentclass[11pt]{article}
\usepackage{amsmath, amssymb, amsthm, bbm}
\usepackage{graphicx}
\usepackage{hyperref}
\usepackage{tikz, subcaption}
\usepackage{enumitem}
\usepackage{geometry}
\usepackage{algorithm}
\usepackage{algpseudocode}
\usepackage[backend=biber,
    maxnames=10,
    minnames=1,
style=alphabetic]{biblatex}
\addbibresource{references.bib}
\geometry{margin=1in}

\newtheorem{proposition}{Proposition}
\newtheorem{theorem}{Theorem}
\newtheorem{lemma}{Lemma}
\newtheorem{corollary}{Corollary}
\newtheorem{definition}{Definition}[section] 

\title{Optimization of Graph Squares - Proofs}
\author{Eddie Xu\\\small\texttt{eddie.xu@student.unimelb.com.au}}
\date{\today}


\begin{document}

\maketitle

\tableofcontents

\section{Formulations} \label{sec::formulations}

Let known constants $g_{ij} \in \{0, 1\}$ come from the adjacency matrix of the square $A_G$ and let $M \in \mathbb{Z}_{> 0}$ be a sufficiently large constant. The variables in the program are $x_{ij}$, $y_{ikj}$, $z_{ij}$ where $x_{ij}$ form the $ij$-th entry of the adjacency matrix of the root $A_H$.  

\subsection{Integer Program} \label{subsec::integer_program}

\begin{alignat}{3}
\text{Variables: } & \nonumber \\ 
& x_{ij} \in \mathbb{Z}_{\geq 0}, 
&& \forall i,j \in \{1, \dots, n\}, \label{eq::IP_x_variable} \\ 
& y_{ikj} \in \mathbb{Z}_{\geq 0}, 
&& \forall i,k,j \in \{1, \dots, n\}, \label{eq::IP_y_variable} \\
& z_{ij} \in \mathbb{Z}_{\geq 0}, 
&& \forall i,j \in \{1, \dots, n\}. \label{eq::IP_z_variable} \\
\text{Constraints: } & \nonumber \\
\text{(1) Indicator argument: } 
& z_{ii} = \sum_{k = 1}^{n}y_{iki} - \sum_{k = 1}^{n} x_{ik}, 
&& \forall i \in \{1, \dots, n\} , \label{eq::IP_fund_matrix} \\
& z_{ij} = \sum_{k = 1}^{n}y_{ikj} + x_{ij}, 
&& \forall i \neq j \in \{1, \dots, n\} , \label{eq::IP_indicator_arg} \\
\text{(2) Indicator function: } 
& z_{ij} \leq M g_{ij}, 
&& \forall i, j \in \{1, \dots, n\}, \label{eq::IP_indicator_upper} \\
& z_{ij} \geq g_{ij}, 
&& \forall i, j \in \{1, \dots, n\}, \label{eq::IP_indicator_lower}\\
\text{(3) Product: } 
& y_{ikj} \leq x_{ik}, 
&& \forall i, k, j \in \{1, \dots, n\}, \label{eq::IP_product_upper_first} \\
& y_{ikj} \leq x_{kj}, 
&& \forall i, k, j \in \{1, \dots, n\}, \label{eq::IP_product_upper_second} \\
& y_{ikj} \geq x_{ik} + x_{kj} - 1, 
&& \forall i, k, j \in \{1, \dots, n\} \label{eq::IP_product_lower} \\
\text{(4) Upper Bound: } 
& x_{ij} \leq 1, 
&& \forall i, j \in \{1, \dots, n\}, \label{eq::IP_x_upper_bound} \\
& y_{ikj} \leq 1, 
&& \forall i, k, j \in \{1, \dots, n\}, \label{eq::IP_y_upper_bound} \\
\text{(5) Symmetry: } 
& x_{ij} = x_{ji}, 
&&\quad \forall i,j \in \{1, \dots, n\}, \label{eq::IP_symmetry_x} \\
\text{(6) Diagonal: } 
& x_{ii} = 0, 
&&\quad \forall i \in \{1, \dots, n\}.  \label{eq::IP_diagonal_x} 
\end{alignat} \label{eq::integer_program}

\subsection{Linear Program} \label{subsec::linear_program}

\begin{alignat}{3}
& x_{ij} \in  \mathbb{R}_{\geq 0}, 
&&\qquad \forall i,j \in \{1, \dots, n\}, \label{eq::LP_x_variable} \\
& y_{ikj} \in \mathbb{R}_{\geq 0}, 
&&\qquad \forall i,k,j \in \{1, \dots, n\}, \label{eq::LP_y_variable} \\
& z_{ij} \in \mathbb{R}_{\geq 0}, 
&&\qquad \forall i,j \in \{1, \dots, n\}, \label{eq::LP_z_variable} \\
\text{Constraints: } & \nonumber \\
\text{(1) Indicator argument: } 
& z_{ii} = \sum_{k = 1}^{n}y_{iki} - \sum_{k = 1}^{n} x_{ik}, 
&& \forall i \in \{1, \dots, n\} , \label{eq::LP_fund_matrix} \\
& z_{ij} = \sum_{k = 1}^{n}y_{ikj} + x_{ij}, 
&& \forall i \neq j \in \{1, \dots, n\} , \label{eq::LP_indicator_arg} \\
\text{(2) Indicator function: } 
& z_{ij} \leq M g_{ij}, 
&& \forall i, j \in \{1, \dots, n\}, \label{eq::LP_indicator_upper} \\
& z_{ij} \geq g_{ij}, 
&& \forall i, j \in \{1, \dots, n\}, \label{eq::LP_indicator_lower}\\
\text{(3) Product: } 
& y_{ikj} \leq x_{ik}, 
&& \forall i, k, j \in \{1, \dots, n\}, \label{eq::LP_product_upper_first} \\
& y_{ikj} \leq x_{kj}, 
&& \forall i, k, j \in \{1, \dots, n\}, \label{eq::LP_product_upper_second} \\
& y_{ikj} \geq x_{ik} + x_{kj} - 1, 
&& \forall i, k, j \in \{1, \dots, n\} \label{eq::LP_product_lower} \\
\text{(4) Upper Bound: } 
& x_{ij} \leq 1, 
&& \forall i, j \in \{1, \dots, n\}, \label{eq::LP_x_upper_bound} \\
& y_{ikj} \leq 1, 
&& \forall i, k, j \in \{1, \dots, n\}, \label{eq::LP_y_upper_bound} \\
\text{(5) Symmetry: } 
& x_{ij} = x_{ji}, 
&&\quad \forall i,j \in \{1, \dots, n\}, \label{eq::LP_symmetry_x} \\
\text{(6) Diagonal: } 
& x_{ii} = 0, 
&&\quad \forall i \in \{1, \dots, n\}.  \label{eq::LP_diagonal_x} 
\end{alignat} \label{eq::linear_program}

\subsection{Max Linear Program} \label{subsec::max_linear_program}

\begin{alignat}{3}
& x_{ij} \in  \mathbb{R}_{\geq 0}, 
&&\qquad \forall i,j \in \{1, \dots, n\}, \label{eq::max_LP_x_variable} \\
& y_{ikj} \in \mathbb{R}_{\geq 0}, 
&&\qquad \forall i,k,j \in \{1, \dots, n\}, \label{eq::max_LP_y_variable} \\
& z_{ij} \in \mathbb{R}_{\geq 0}, 
&&\qquad \forall i,j \in \{1, \dots, n\}, \label{eq::max_LP_z_variable} \\
\text{Constraints: } & \nonumber \\
\text{(1) Indicator argument: } 
& z_{ii} = \sum_{k = 1}^{n}y_{iki} - \sum_{k = 1}^{n} x_{ik}, 
&& \forall i \in \{1, \dots, n\} , \label{eq::max_LP_fund_matrix} \\
& z_{ij} = \sum_{k = 1}^{n}y_{ikj} + x_{ij}, 
&& \forall i \neq j \in \{1, \dots, n\} , \label{eq::max_LP_indicator_arg} \\
\text{(2) Indicator function: } 
& z_{ij} \leq M g_{ij}, 
&& \forall i, j \in \{1, \dots, n\}, \label{eq::max_LP_indicator_upper} \\
& z_{ij} \geq g_{ij}, 
&& \forall i, j \in \{1, \dots, n\}, \label{eq::max_LP_indicator_lower}\\
\text{(3) Product: } 
& y_{ikj} \leq x_{ik}, 
&& \forall i, k, j \in \{1, \dots, n\}, \label{eq::max_LP_product_upper_first} \\
& y_{ikj} \leq x_{kj}, 
&& \forall i, k, j \in \{1, \dots, n\}, \label{eq::max_LP_product_upper_second} \\
& y_{ikj} \geq x_{ik} + x_{kj} - 1, 
&& \forall i, k, j \in \{1, \dots, n\} \label{eq::max_LP_product_lower} \\
\text{(4) Upper Bound: } 
& x_{ij} \leq 1, 
&& \forall i, j \in \{1, \dots, n\}, \label{eq::max_LP_x_upper_bound} \\
& y_{ikj} \leq 1, 
&& \forall i, k, j \in \{1, \dots, n\}, \label{eq::max_LP_y_upper_bound} \\
\text{(5) Symmetry: } 
& x_{ij} = x_{ji}, 
&&\quad \forall i,j \in \{1, \dots, n\}, \label{eq::max_LP_symmetry_x} \\
\text{(6) Diagonal: } 
& x_{ii} = 0, 
&&\quad \forall i \in \{1, \dots, n\},  \label{eq::max_LP_diagonal_x} \\
\text{Objective Function: } & \nonumber \\ 
& \max \sum_{i=1}^{n} \sum_{j=1}^{n} x_{ij}. \label{eq::max_objective_funtion}
\end{alignat} \label{eq::max_linear_program}

\section{Two-Edge-Triangle Proof} \label{sec::two_edge_triangle_proof}

Let $x$ be an optimal solution to the linear relaxation defined by Equations~\eqref{eq::max_LP_x_variable}--\eqref{eq::max_objective_funtion}. 
If $g_{ij}=0$, then for every $k\in\{1,\dots,n\}$,
\[
x_{ik}+x_{kj}=1.
\]

\subsection{Integer Program} \label{subsec::triangle_two_edges_integer_program}

\begin{lemma} \label{lemma::two_edges_integer_sum}
  Consider the triangle with three vertices $i$, $j$ and $k$ where edges $ik$ and $kj$ are present and $ij$ is absent in square $G$. Then for root $G$, the integer program defined by Equations~\eqref{eq::IP_x_variable}-\eqref{eq::IP_diagonal_x} must satisfy
  \[
  x_{ik} + x_{kj} \leq 1. 
  \]
\end{lemma}
\begin{proof}
  As $ij$ is absent in $G$, $g_{ij} = 0$ such that Equation~\eqref{eq::IP_indicator_upper} gives $z_{ij} \leq 0$ and Equation~\eqref{eq::IP_indicator_lower} gives $z_{ij} \geq 0$ such that $z_{ij} = 0$. Therefore Equation~\eqref{eq::IP_indicator_arg} gives 
  \[
  0 = \sum_{a=1}^{n}y_{iaj}+x_{ij} = \sum_{a=1}^{n}x_{ia}x_{aj}+x_{ij},
  \]
  where the second equality follows from the product constraints Equations~\eqref{eq::IP_product_upper_first}--\eqref{eq::IP_product_lower} together with the integrality of $x \in \{0, 1\}$ by Equation~\eqref{eq::IP_x_variable} and Equation~\eqref{eq::IP_x_upper_bound}. Since all variables are non-negative, $0$ only holds when $x_ij = 0$ and all of$x_{ia} x_{aj}=0$ and using the integrality of $x$, this means that at most one of $x_{ia}$, $x_{aj}$ is equal to $1$ and the other must be $0$, which can be equivalently written as $x_{ia} + x_{aj} \leq 1$. Taking $a=k$ gives $x_{ik}+x_{kj} \leq 1$ as required. 
\end{proof}

\subsection{Linear Program} \label{subsec::triangle_two_edges_linear_program}

\begin{lemma} \label{lemma::two_edges_linear_sum}
  Consider the triangle with three vertices $i$, $j$ and $k$ where edges $ik$ and $kj$ are present and $ij$ is absent in square $G$. Then for root $G$, the linear program defined by Equations~\eqref{eq::LP_x_variable}-\eqref{eq::LP_diagonal_x} must satisfy 
  \[
  x_{ik} + x_{kj} \leq 1. 
  \]
\end{lemma}
\begin{proof}
  Since $g_{ij}=0$, Equations~\eqref{eq::LP_indicator_upper}--\eqref{eq::LP_indicator_lower} imply that $z_{ij}=0$.
  Equation~\eqref{eq::LP_indicator_arg} therefore gives
  \[
  0 = z_{ij} = \sum_{a=1}^{n} y_{iaj} + x_{ij}.
  \]
  Since all variables are nonnegative by Equations~\eqref{eq::LP_x_variable}--\eqref{eq::LP_z_variable}, every term in the above sum must be equal to zero with $x_{ij} = 0$ and $y_{iaj} = 0$ for all $a$. Thus applying Equation~\eqref{eq::LP_product_lower} to $y_{ikj} = 0$ gives
  \[
  0 = y_{ikj} \geq x_{ik} + x_{kj} - 1,
  \]
  and hence $x_{ik}+x_{kj} \leq 1$ as required.
\end{proof}


% z_{ij} with sum of y_{ikj} must be 0 - however that doesnt necessarily mean all y_{ikj} are 0 - there could be one that is negative since y_{ikj} >= x_{ik} + x_{kj} - 1 which means 
% the lower bound of y_{ikj} could be negative when $x_{ik} + x_{kj} < 1$. if however it is negative,  that means that another y_{ikj} must be positive OR x_{ij} must be positive.

\subsection{Max Linear Program} \label{subsec::triangle_two_edges_max_linear_program}

\begin{proposition} \label{prop::local_two_edges_augmentation}
  Consider the maximum linear relaxation defined by Equations~\eqref{eq::max_LP_x_variable}--\eqref{eq::max_objective_funtion} on the graph induced by the three vertices $\{i,j,k\}$. Suppose
  \[
  g_{ik}=1,\quad g_{kj}=1,\quad g_{ij}=0,
  \]
  then there always exists an optimal, feasible solution satisfying $x_{ik}+x_{kj}=1$.
\end{proposition}
\begin{proof}
  Since $g_{ij}=0$, Equations~\eqref{eq::max_LP_indicator_upper}--\eqref{eq::max_LP_indicator_lower} thus imply that $z_{ij}=0$. Equation~\eqref{eq::max_LP_indicator_arg} therefore gives
  \begin{equation} \label{eq::local_two_edges_augementation_indicator_arg}
    0 = \sum_{t\in\{i,j,k\}} y_{itj} + x_{ij} = y_{iij} + y_{ijj} + y_{ikj} + x_{ij} = y_{ikj} + x_{ij},
  \end{equation}
  since Equations~\eqref{eq::max_LP_product_upper_first}--\eqref{eq::max_LP_product_upper_second} together with Equation~\eqref{eq::max_LP_diagonal_x} give that $y_{iij} \leq x_{ii} = 0$ and $y_{ijj} \leq x_{jj} = 0$ and as all variables are non-negative by Equations~\eqref{eq::max_LP_x_variable}--\eqref{eq::max_LP_z_variable}, it follows that $y_{iij} = 0$, $y_{ijj} = 0$. Therefore Equation~\eqref{eq::local_two_edges_augementation_indicator_arg} gives
  \[
  y_{ikj} = 0, \quad x_{ij}=0. 
  \]

  By contradiction, assume that $x_{ik}$ and $x_{kj}$ give a optimal, feasible solution to the maximal linear program where $x_{ik} + x_{kj} = 1 - \delta$ and $0 \leq \delta \leq \varepsilon$ such that $x_{ik} + x_{kj} < 1$. Consider the perturbation $x_{ik}' = x_{ik} + \epsilon$ such that $x_{ik}' + x_{kj} = 1$. 
\end{proof}

\begin{proposition} \label{prop::extended_two_edges_augmentation}
  Consider the maximum linear relaxation defined by Equations~\eqref{eq::max_LP_x_variable}--\eqref{eq::max_objective_funtion} on the square graph $G$. Suppose
  \[
  g_{ik}=1,\quad g_{kj}=1,\quad g_{ij}=0,
  \]
  and let $(x,y,z)$ be a feasible solution satisfying $x_{ik}+x_{kj}<1$. Then in order for there to exist another feasible solution $(x',y',z')$ such that $ \displaystyle \sum_{u,v\in\{i,j,k\}}x'_{uv} > \sum_{u,v\in\{i,j,k\}}x_{uv}$ $G$ must be infinitely large.
\end{proposition}
\begin{proof}
  Using Proposition~\ref{prop::local_two_edges_augmentation}, any three vertices that satisfy the maximum linear relaxation defined by Equations~\eqref{eq::max_LP_x_variable}--\eqref{eq::max_objective_funtion} will always satisfy $x_{ik}+x_{kj}=1$. Hence the only way $x_{ik} + x_{kj} < 1$ is if 
\end{proof}

\begin{lemma} \label{prop::two_edges_augmentation}
  Let $x$ be a feasible solution to the maximum linear relaxation defined by Equations~\eqref{eq::max_LP_x_variable}--\eqref{eq::max_objective_funtion}. Suppose there exists vertices $i$, $j$ and $k$ such that
  \[
  g_{ik}=1,\quad g_{kj}=1,\quad g_{ij}=0 \quad \text{and} \quad x_{ik} + x_{kj} < 1. 
  \]
  Then there exists another feasible solution $(x',y',z')$ satisfying $\displaystyle \sum_{u,v} x'_{uv} > \sum_{u,v} x_{uv}$.
\end{lemma}
\begin{proof}

\end{proof}

\begin{theorem} \label{theorem::two_edges_linear_maximization}
  Let $x$ be an optimal solution to the linear relaxation defined by Equations~\eqref{eq::max_LP_x_variable}--\eqref{eq::max_objective_funtion}.
  If $g_{ij}=0$, then for every $k\in\{1,\dots,n\}$,
  \[
  x_{ik}+x_{kj}=1.
  \]
\end{theorem}
\begin{proof}
  Using Lemma~\ref{lemma::two_edges_linear_sum} $x_{ik} + x_{kj} \leq 1$ and using Proposition~\ref{prop::two_edges_augmentation} $x_{ik} + x_{kj} \geq 1$
  such that $x_{ik}+x_{kj}=1$ as required. 
\end{proof}





\end{document}
